\section{Introduction}\label{sec:intro}

\texdocx{} is a compiler that turns a \LaTeX{} document into a Word file
(\texttt{.docx}). It does not run \TeX{} and it does not print to PDF first.
Instead it rebuilds the part of \LaTeX{} that \emph{reads} your source---the
tokenizer with its category codes, macro expansion, document structure and
mathematics---and hands the result to Word, which does the layout itself when
it opens the file. Line breaking, page breaking and font metrics are Word's job,
so the output reflows like any document typed in Word.

Three properties set \texdocx{} apart from converters that only approximate the
look of a \LaTeX{} paper:
\begin{itemize}
  \item \textbf{Equations stay editable.} Every formula becomes native Office
        Math (OMML), so a co-author can click into it and change a subscript in
        Word's equation editor. Section~\ref{sec:math} covers the details.
  \item \textbf{Numbers are live.} References to sections, equations,
        figures and tables, the equation and caption numbers themselves and
        the table of contents are written as Word fields with the correct value
        already filled in, so the document looks right as soon as it opens and
        Word can update the numbers after you edit it.
  \item \textbf{Your macros work.} \texdocx{} has a real macro expander, so
        \cmd{newcommand}, \cmd{def}, conditionals and environments you define
        yourself are expanded by \TeX's own rules instead of being
        matched against patterns (Section~\ref{sec:macros}).
\end{itemize}

This guide is itself a test of those claims: it was written in \LaTeX{} and
converted by \texdocx{}. Section~\ref{sec:features} surveys the supported
document features, Section~\ref{sec:figures} covers figures and tables, and
Section~\ref{sec:security} covers safety and diagnostics, including the sandbox
that keeps a document from reading anything outside its own folder.

\subsection{Quick start}

\texdocx{} needs Node.js~23.6 or newer and the pnpm package
manager.\footnote{The compiler is written in TypeScript and runs directly
through Node's built-in type stripping, which is why there is no build step and
why older Node versions are not supported.} From the repository root, install
the dependencies once and then convert a file:
\begin{lstlisting}[language=bash]
pnpm install
pnpm texdocx paper.tex                 # writes paper.docx next to it
pnpm texdocx paper.tex -o out.docx --json --strict
\end{lstlisting}

Problems are reported the way a compiler reports them: a severity, a code such
as \texttt{W013}, the file, line and column, the offending source line and,
where possible, a hint. The command accepts the following options:
\begin{description}
  \item[\texttt{-o}, \texttt{-{}-output} \textit{file}] Where to write the
        Word file. By default the output takes the input's name with the
        extension \texttt{.docx}.
  \item[\texttt{-{}-root} \textit{dir}] The folder the document may read
        from. By default this is the folder that contains the input file;
        \cmd{input} and \cmd{includegraphics} cannot reach outside it.
  \item[\texttt{-{}-strict}] Treat every warning as an error, for example to
        make a build fail on an unknown command.
  \item[\texttt{-{}-json}] Print each diagnostic as one line of JSON instead of
        the human-readable form, for editors and scripts.
  \item[\texttt{-q}, \texttt{-{}-quiet}] Print errors only.
  \item[\texttt{-h}, \texttt{-{}-help}] Show a short usage summary.
\end{description}
The exit status is~0 on success, 1 when the document had errors and 2 when the
command line was wrong or the input file could not be read. The Word file is written even when there are
errors, so you can open it and see what went wrong.

\begin{tip}
Run \texdocx{} with \texttt{-{}-strict} in continuous integration. A warning
about an unknown command usually means some text in the Word file is not what
you meant.
\end{tip}

\subsection{How it works}

\texdocx{} is designed as a pipeline of six stages. Only the first three
resemble \TeX; the last three replace \TeX's typesetter with a typed document
model and a Word writer. As in \TeX, the first four stages run interleaved,
one token at a time: the digester asks the expander for the next token, the
expander asks the lexer, and a file is loaded only when the command that
names it runs.
\begin{enumerate}
  \item \textbf{Load.} Every file the document pulls in through
        \cmd{input}, \cmd{include} or \cmd{includegraphics} is read through
        a sandbox confined to the project folder, and every token remembers
        where it came from, so a diagnostic can name the file, line and column.
  \item \textbf{Tokenize.} The lexer turns characters into tokens, looking
        up each character's category code as it reads it, so a
        \cmd{makeatletter} or \cmd{catcode} in the document changes how the
        very next character is read.
  \item \textbf{Expand.} The expander replaces macros by their definitions and
        evaluates conditionals and groups, within budgets on steps, nesting
        depth, token count and time so that a runaway macro stops with an error.
  \item \textbf{Digest.} Command and environment handlers build a document
        model of headings, paragraphs, lists, tables, figures and math trees
        instead of \TeX's boxes and glue.
  \item \textbf{Resolve.} Once the whole document has been read, every
        \cmd{ref} is matched to its \cmd{label}, so forward references work
        in a single run where \LaTeX{} needs two. An undefined reference
        prints \texttt{??} and a warning, as in \LaTeX.
  \item \textbf{Write.} The writer turns the model into Word's XML---styles,
        numbering, fields, footnotes, images and equations---and packs it into
        a byte-for-byte reproducible \texttt{.docx} archive.
\end{enumerate}
The writer depends only on the document model and never sees \LaTeX{} source,
which keeps the quirks of each format on its own side of the model.
